\documentclass[11pt]{article}
\usepackage[margin=1in]{geometry}
\usepackage{enumitem}
% =============================================================================
% Preambles/header.tex — shared LaTeX/Beamer preamble for this project
%
% Use from a Beamer lecture by prepending:
%     \input{header}
% and compiling with TEXINPUTS=../Preambles:$TEXINPUTS (the /compile-latex
% skill does this automatically).
%
% PALETTE CONTRACT. Color names defined here MUST match the names in
% Quarto/theme-template.scss so Beamer and Quarto renderings use the same
% palette. The scripts/check-palette-sync.sh script enforces this.
%
% To customize: edit the HEX values below AND the matching $variable in
% Quarto/theme-template.scss, then run scripts/check-palette-sync.sh.
% =============================================================================

% -----------------------------------------------------------------------------
% Engine detection + encoding
%
% Our /compile-latex skill uses XeLaTeX, where inputenc is unnecessary and
% can produce encoding warnings. Only load inputenc under pdfTeX.
% -----------------------------------------------------------------------------
\usepackage{iftex}
\ifPDFTeX
  \usepackage[utf8]{inputenc}
\fi

\usepackage{amsmath, amssymb, amsthm}
\usepackage{booktabs}
\usepackage{graphicx}

% -----------------------------------------------------------------------------
% PALETTE — mirrors Quarto/theme-template.scss variable names.
% Load xcolor and define colors BEFORE anything that references them
% (notably hyperref, hypersetup, and the Beamer color assignments).
% -----------------------------------------------------------------------------
\usepackage{xcolor}

\definecolor{primary-blue}{HTML}{012169}      % headings, accents
\definecolor{primary-gold}{HTML}{B9975B}      % emphasis, borders
\definecolor{highlight-yellow}{HTML}{F2A900}  % markers, alerts
\definecolor{light-bg}{HTML}{E8EDF5}          % title / section backgrounds
\definecolor{jet}{HTML}{1A1A1A}               % body text

% Semantic colors (used in diagrams and annotations)
\definecolor{positive}{HTML}{15803D}          % good / observed / identified
\definecolor{negative}{HTML}{B91C1C}          % bad / problematic / bias
\definecolor{neutral}{HTML}{525252}           % reference / context

% Highlight accents (match .hi-* classes in the SCSS)
\definecolor{hi-slate}{HTML}{314F4F}
\definecolor{hi-green}{HTML}{2E8B57}
\definecolor{hi-red}{HTML}{C41E3A}

% Semantic aliases so skill templates and snippets can write
% \textcolor{accent}{...} without hard-coding "primary-blue".
\colorlet{accent}{primary-blue}
\colorlet{accent2}{primary-gold}

% -----------------------------------------------------------------------------
% Hyperref — loaded AFTER xcolor + palette so link colors resolve.
% -----------------------------------------------------------------------------
\usepackage{hyperref}
\hypersetup{colorlinks=true, linkcolor=primary-blue, urlcolor=primary-blue,
            citecolor=primary-blue, pdfborder={0 0 0}}

% -----------------------------------------------------------------------------
% Course code — set once per deck (after \input{header}, before \begin{document})
% via \coursecode{CS401}. Shown in the footer on every slide so a deck's course
% is identifiable without opening the title page. Empty by default.
% -----------------------------------------------------------------------------
\makeatletter
\newcommand{\coursecode}[1]{\gdef\@coursecodetext{#1}}
\gdef\@coursecodetext{}
\makeatother

% -----------------------------------------------------------------------------
% Beamer theme assignments (only apply under Beamer).
%
% We detect Beamer via \@ifclassloaded{beamer}, which is the documented,
% non-internal way. This must be wrapped in \makeatletter/\makeatother
% because \@ifclassloaded contains an @-catcode character.
% -----------------------------------------------------------------------------
\makeatletter
\@ifclassloaded{beamer}{%
  \usefonttheme{professionalfonts}
  \setbeamercolor{structure}{fg=primary-blue}
  \setbeamercolor{frametitle}{fg=primary-blue}
  \setbeamercolor{title}{fg=primary-blue}
  \setbeamercolor{subtitle}{fg=primary-gold}
  \setbeamercolor{author}{fg=primary-blue}
  \setbeamercolor{institute}{fg=neutral}
  \setbeamercolor{date}{fg=primary-gold}
  \setbeamercolor{itemize item}{fg=primary-blue}
  \setbeamercolor{itemize subitem}{fg=primary-gold}
  \setbeamercolor{alerted text}{fg=primary-gold}
  \setbeamercolor{block title}{fg=white, bg=primary-blue}
  \setbeamercolor{block body}{bg=light-bg}
  % Semantic block variants: block=definition/concept (blue), exampleblock=
  % worked example (green/positive), alertblock=key takeaway or caution (gold).
  % Keep to <=2 colored boxes per slide across ALL three variants (INV-7).
  \setbeamercolor{block title example}{fg=white, bg=positive}
  \setbeamercolor{block body example}{bg=light-bg}
  \setbeamercolor{block title alerted}{fg=white, bg=primary-gold}
  \setbeamercolor{block body alerted}{bg=light-bg}

  % Minimal footer: course code (left) + page X / N (right), no navigation clutter
  \setbeamertemplate{navigation symbols}{}
  \setbeamertemplate{footline}{%
    \color{neutral}\footnotesize\hspace{0.7em}\@coursecodetext\hfill
    \insertframenumber/\inserttotalframenumber\hspace{0.7em}\vspace{0.3em}}
}{}
\makeatother

% -----------------------------------------------------------------------------
% TikZ setup — libraries the snippet gallery uses
% -----------------------------------------------------------------------------
\usepackage{tikz}
\usetikzlibrary{arrows.meta, positioning, calc, decorations.pathreplacing,
                fit, shapes.geometric, backgrounds}

% Shared TikZ styles — snippets and hand-written diagrams can pick these up
% via `\begin{tikzpicture}[every node/.style={...}]` or by naming specific
% styles (e.g., `\node[dag-node] (X) at (0,0) {X};`).
\tikzset{
  % Node styles for causal DAGs and similar
  dag-node/.style = {
    draw=primary-blue, fill=light-bg, rounded corners=2pt,
    minimum width=1.2cm, minimum height=0.9cm, align=center, font=\small
  },
  decision-node/.style = {
    draw=primary-gold, fill=white, diamond, aspect=1.8,
    minimum width=1.8cm, minimum height=1.2cm, align=center, font=\small,
    inner sep=1pt
  },
  % Edge styles — observed vs counterfactual
  observed-edge/.style = {-{Latex[length=2.5mm]}, thick, draw=jet},
  counterfactual-edge/.style = {-{Latex[length=2.5mm]}, thick, draw=neutral, dashed},
  confound-edge/.style = {-{Latex[length=2.5mm]}, thick, draw=negative, dashed},
  % Data-point styles for plots
  observed-dot/.style = {fill=primary-blue, draw=primary-blue, circle, inner sep=1.2pt},
  counterfactual-dot/.style = {fill=white, draw=neutral, circle, inner sep=1.2pt}
}

% -----------------------------------------------------------------------------
% Convenience macros
% -----------------------------------------------------------------------------
\newcommand{\muted}[1]{\textcolor{neutral}{#1}}
\newcommand{\key}[1]{\textcolor{primary-gold}{\textbf{#1}}}
\newcommand{\good}[1]{\textcolor{positive}{#1}}
\newcommand{\bad}[1]{\textcolor{negative}{#1}}

% Standout frame macro: full-bleed dark background for section transitions.
% Beamer-only (uses \begin{frame}/\setbeamercolor/\usebeamerfont) -- do not
% call from an article-class file (e.g. Notes/<CODE>/*-notes.tex). \muted,
% \key, \good, \bad, and \coursecode above are plain xcolor/gdef and are
% safe to use from either class; verified when Notes/article-class support
% was added.
% Expand: \transitionslide{Your transition title}
\newcommand{\transitionslide}[1]{%
  \begingroup
  \setbeamercolor{background canvas}{bg=primary-blue}
  \begin{frame}[plain]
    \centering
    \vfill
    {\usebeamerfont{title}\color{white}\Large #1\par}
    \vfill
  \end{frame}
  \endgroup
}

% Section divider frame (Beamer-only), an alternative to \transitionslide with a
% darker two-line style: near-black (jet) background, a small white label line
% above a larger gold title, and a thin gold rule along the bottom edge.
% Expand: \sectiondivider{Section 2}{Pointers}
\newcommand{\sectiondivider}[2]{%
  \begingroup
  \setbeamercolor{background canvas}{bg=jet}
  \begin{frame}[plain]
    \vspace*{3.0cm}
    {\color{white}\ttfamily\large #1\par}
    \vspace{0.5cm}
    {\color{primary-gold}\LARGE #2\par}
    \begin{tikzpicture}[remember picture, overlay]
      \draw[primary-gold, line width=2.5pt]
        ($(current page.south west)+(0,0.1)$) -- ($(current page.south east)+(0,0.1)$);
    \end{tikzpicture}
  \end{frame}
  \endgroup
}

% -----------------------------------------------------------------------------
% Channel watermark — "Concepts That Click" (YouTube).
%
% Article-class files (CompetitiveExam Q/A sets, Notes, Assignments) get a
% light diagonal draft watermark on every page plus a centered footer naming
% the channel. Beamer decks get a subtle TikZ background watermark instead
% (the footline already carries the course code). Centralized here so every
% MCQ PDF is branded without per-file edits.
% -----------------------------------------------------------------------------
\makeatletter
\@ifclassloaded{article}{%
  \usepackage{draftwatermark}
  \SetWatermarkText{Concepts That Click}
  \SetWatermarkScale{1}
  \SetWatermarkFontSize{2cm}
  \SetWatermarkLightness{0.86}
  \SetWatermarkAngle{45}
  \usepackage{fancyhdr}
  \pagestyle{fancy}
  \fancyhf{}
  \fancyfoot[C]{\footnotesize\textcolor{neutral}{Concepts That Click~|~YouTube: \href{https://www.youtube.com/@concepts.thatclick}{@concepts.thatclick}}}
  \renewcommand{\headrulewidth}{0pt}
  \renewcommand{\footrulewidth}{0pt}
  \fancypagestyle{plain}{%
    \fancyhf{}%
    \fancyfoot[C]{\footnotesize\textcolor{neutral}{Concepts That Click~|~YouTube: \href{https://www.youtube.com/@concepts.thatclick}{@concepts.thatclick}}}%
    \renewcommand{\headrulewidth}{0pt}%
    \renewcommand{\footrulewidth}{0pt}%
  }
}{}
\@ifclassloaded{beamer}{%
  \setbeamertemplate{background}{%
    \begin{tikzpicture}[remember picture, overlay]
      \node[rotate=45, opacity=0.055, align=center]
        at (current page.center)
        {\Huge\textcolor{neutral}{Concepts That Click}};
    \end{tikzpicture}%
  }
}{}
\makeatother 
\coursecode{JKSSB}
\renewcommand{\thesection}{33.\arabic{section}}
\newtheorem{example}{Example}[section]
\newenvironment{definitionbox}[1]{%
  \par\noindent\textbf{Definition (#1).}\ \itshape
}{\par\vspace{0.3em}}
\title{PDF and File Formats --- Detailed Lecture Notes}
\author{JKSSB Computer Awareness}
\date{\today}

\begin{document}
\maketitle

\section{What a file format is}
Every file on a computer is stored as a long series of bits. On its own, a
string of bits says nothing about whether it is a letter, a photograph, a
spreadsheet, or a slide show. The \textbf{file format} is the rule that tells
the software how those bits are arranged --- where the text begins, where the
image data sits, and how the page should be drawn. When a program opens a
file, it reads the format to know how to interpret the contents.

The \textbf{extension} is the short suffix at the end of a file name, such as
\texttt{.pdf}, \texttt{.docx}, or \texttt{.png}. The extension is a label
that helps the operating system choose which program should open the file.
It is convenient but only a name: the real identity of the file is its format.

\begin{definitionbox}{File format}
The defined structure or rule by which data is stored in a file, so that
software can correctly read, display, and save its contents.
\end{definitionbox}

\begin{definitionbox}{Extension}
The suffix at the end of a file name, usually written with a dot, that
indicates the file's intended type (for example \texttt{.pdf}, \texttt{.docx},
\texttt{.png}).
\end{definitionbox}

\begin{center}
\begin{tabular}{p{0.28\textwidth}p{0.58\textwidth}}
\toprule
\textbf{Term} & \textbf{Meaning in this lesson} \\
\midrule
Format & The actual structure of the data inside the file. \\
Extension & The dot-suffix in the file name, such as \texttt{.docx}. \\
Fixed layout & A format that stores exact positions, so the page looks the same everywhere. \\
Plain text & A file of ordinary characters with no formatting. \\
\bottomrule
\end{tabular}
\end{center}

\section{Extension versus format (the core trap)}
The single most repeated confusion in this topic is treating the extension and
the format as the same thing. They are not. In the file name
\texttt{report.docx}, the base name is \texttt{report}, the extension is
\texttt{docx} (written \texttt{.docx}), and the format is DOCX, the Office
Open XML word-processing format. Likewise, \texttt{.png} is an extension,
while PNG (Portable Network Graphics) is the format.

The extension can also be changed easily, and doing so does not convert the
file. Renaming a photograph from \texttt{pic.jpg} to \texttt{pic.png} does
not turn a JPEG image into a PNG image; it only changes the label. This is
why the exam stresses the distinction: the extension is a name, the format is
the content.

\begin{example}
An item states, ``\texttt{.png} is an image file format.'' The statement
mixes the two ideas. \texttt{.png} is the \textbf{extension}; PNG is the
\textbf{format}. A correct statement would be, ``PNG is an image file format,
and \texttt{.png} is its usual extension.''
\end{example}

\section{Common document formats}
The Office applications each save their work in a standard format. These are
the Office Open XML formats, and each has its own extension.

\begin{center}
\begin{tabular}{p{0.16\textwidth}p{0.30\textwidth}p{0.40\textwidth}}
\toprule
\textbf{Extension} & \textbf{Format} & \textbf{Typical use} \\
\midrule
\texttt{.docx} & Word document & Text, formatting, and images (MS Word) \\
\texttt{.xlsx} & Excel workbook & Worksheets, formulas, and charts (MS Excel) \\
\texttt{.pptx} & PowerPoint presentation & Slides and slide shows (PowerPoint) \\
\texttt{.csv} & Comma-Separated Values & Plain-text rows with values separated by commas \\
\texttt{.txt} & Plain text & Unformatted characters only \\
\bottomrule
\end{tabular}
\end{center}

A \textbf{word-processing document} (DOCX) stores formatted text. A
\textbf{workbook} (XLSX) stores worksheets with formulas and numbers. A
\textbf{presentation} (PPTX) stores slides. A \textbf{CSV} file is plain text
in which each line is a record and commas separate the fields; it is often
used to move data between programs. A \textbf{TXT} file holds plain
characters and no formatting at all.

\section{Common image extensions}
Images also come in many formats, each suited to a different purpose.
\textbf{PNG} (Portable Network Graphics) uses lossless compression and
supports transparency, so it is good for logos and screenshots.
\textbf{JPG/JPEG} uses lossy compression, so it is good for photographs where
a small loss of quality is acceptable. \textbf{GIF} (Graphics Interchange
Format) handles simple images and short animations. \textbf{BMP} (Bitmap)
stores a raw image and is usually larger. \textbf{TIFF} (Tagged Image File
Format) is used for high-quality scanning and printing. \textbf{SVG}
(Scalable Vector Graphics) stores a drawing as lines and shapes, so it scales
without blurring.

\begin{center}
\begin{tabular}{p{0.16\textwidth}p{0.34\textwidth}p{0.36\textwidth}}
\toprule
\textbf{Extension} & \textbf{Format} & \textbf{Character} \\
\midrule
\texttt{.png} & Portable Network Graphics & Lossless, transparency \\
\texttt{.jpg} & JPEG image & Lossy, photos \\
\texttt{.gif} & Graphics Interchange Format & Simple images, animation \\
\texttt{.bmp} & Bitmap & Raw, larger \\
\texttt{.tiff} & Tagged Image File Format & High quality, scanning/print \\
\texttt{.svg} & Scalable Vector Graphics & Vector, scales without blur \\
\bottomrule
\end{tabular}
\end{center}

\section{PDF: the Portable Document Format}
\textbf{PDF} stands for \textbf{Portable Document Format}. It is a
\textbf{fixed-layout} format: every item on the page is placed at a fixed
position. Because the file also carries the fonts and images it needs, a PDF
looks the \textbf{same on every device and printer}. This is why PDF is used
for e-books, forms, official letters, admit cards, and result sheets, where
the exact appearance matters.

Contrast this with a word-processing document, which is a \emph{re-flowable}
format: text re-flows to fit the screen, so the layout can shift when the
font or device changes. A PDF, by design, does not re-flow; it is a picture of
a finished page expressed as instructions.

\begin{definitionbox}{PDF}
Portable Document Format --- a fixed-layout file format that preserves the
appearance of a document so it renders identically across devices and
printers.
\end{definitionbox}

\section{Text and scanned PDFs}
Not every PDF contains real text. It is important to separate two kinds.

A \textbf{text PDF} (sometimes called an editable-source PDF) is generated
from a source document such as a Word file. Its pages contain genuine text and
vector graphics. The text can be selected, copied, and searched, and the file
is usually smaller.

A \textbf{scanned PDF} (image-only PDF) is produced by scanning printed pages
or photographing them. Each page is stored as an \emph{image}, much like a
photograph. The text is readable to the human eye but not to the computer: it
cannot be selected, copied, or searched. Scanned files are also usually
larger.

\begin{example}
An item says, ``In a scanned PDF, you can search for a word using
\texttt{Ctrl+F}.'' This is wrong. A scanned PDF is only a picture of the
page, so its text is not searchable until OCR is applied. In a text PDF, the
same search works immediately.
\end{example}

\section{OCR: Optical Character Recognition}
\textbf{OCR} stands for \textbf{Optical Character Recognition}. It is the
technique that reads the \emph{picture of the text} and converts it into
\textbf{machine-readable text}. Once OCR has been run, the text in a scanned
document can be selected, searched, copied, and edited just like the text in a
text PDF.

OCR is used on scanned pages and on photographs of documents. It is the
bridge that turns an image-only PDF into a searchable one. The recognition is
not always perfect, so the result should be checked.

\begin{definitionbox}{OCR (Optical Character Recognition)}
The process of converting an image of text into machine-readable, editable,
and searchable text.
\end{definitionbox}

\section{Encryption and password protection}
A PDF can be protected so that its contents are confidential.
\textbf{Encryption} scrambles the file so that it cannot be read without the
correct key or password. Two kinds of password are common. A \textbf{password
to open} (user password) is required before the file will open at all. A
\textbf{permissions password} (owner password) restricts what a reader may
do --- for example, printing, copying, or editing.

Encryption provides \textbf{confidentiality}: only the authorised reader can
see the contents. It does not prove who created the file, and it is not the
same as a digital signature.

\section{Digital signature}
A \textbf{digital signature} is attached to a file to prove \textbf{who
signed it}. It also proves \textbf{integrity}: if the file is altered after
it has been signed, the signature no longer matches and the change is
detected. A digital signature uses a \textbf{digital certificate} issued by a
trusted authority.

A digital signature provides \textbf{authenticity} (the signer is verified)
and \textbf{integrity} (the file has not changed). It does \emph{not} hide
the contents. This is the distinction the exam most often tests: a digital
signature is about \emph{who} and \emph{whether it changed}, while encryption
is about \emph{who can read it}.

\begin{definitionbox}{Digital signature}
An electronic signature that verifies the identity of the signer and detects
any change made to the file after signing, providing authenticity and
integrity.
\end{definitionbox}

\begin{center}
\begin{tabular}{lll}
\toprule
\textbf{Point} & \textbf{Encryption} & \textbf{Digital signature} \\
\midrule
Main purpose & Confidentiality & Authenticity and integrity \\
Result & Contents unreadable & Signer verified, tampering detected \\
Needs & Key / password & Digital certificate \\
\bottomrule
\end{tabular}
\end{center}

\begin{example}
An item states, ``A digital signature encrypts the PDF so that only the
recipient can read it.'' The statement is wrong. The encryption that keeps
others from reading the file is \textbf{encryption}. A digital signature does
not hide the content; it confirms the signer and shows whether the file has
been changed.
\end{example}

\section{Redaction and PDF/A}
\textbf{Redaction} means \textbf{permanently removing} sensitive content
before a file is shared. Names, account numbers, addresses, or photographs
may be redacted. A proper redaction deletes the underlying text or image; it
is not enough to draw a black rectangle over it, because hidden text can
often still be selected and copied. A careful reader should always check that
the removed text is truly gone.

\textbf{PDF/A} is a special ISO-standard version of PDF designed for
\textbf{long-term archiving}. It embeds the fonts and colour information
inside the file so that the document will render exactly the same years later,
even on future software. The ``A'' in PDF/A stands for \emph{archival}. It is
used by libraries, courts, and government records, and is meant for
preservation rather than everyday editing.

\section{Exam retrieval and common confusions}
Five distinctions prevent most errors on this topic.
\begin{enumerate}
  \item The \textbf{extension} is the file-name suffix; the \textbf{format}
    is the structure of the data. \texttt{.png} is an extension; PNG is the
    format (TRAP-34).
  \item A \textbf{scanned PDF} is an image; its text cannot be selected or
    searched until \textbf{OCR} is applied. A text PDF is searchable at once.
  \item \textbf{Encryption} gives confidentiality; a \textbf{digital
    signature} gives authenticity and integrity. They are not the same
    (TRM-11, TRAP-15).
  \item \textbf{Redaction} removes sensitive content permanently; a black
    box drawn over text is not sufficient.
  \item \textbf{PDF/A} is for long-term \textbf{archiving}; ordinary PDF is
    the everyday fixed-layout format.
\end{enumerate}

\begin{example}
An item asks which format keeps a fixed layout and looks the same on every
device. The answer is PDF. If the stem instead asks which format stores
spreadsheet data, the answer is XLSX; if it asks which is a plain-text data
format, the answer is CSV; and if it asks which format is meant for long-term
archiving, the answer is PDF/A.
\end{example}

\subsection*{Exercises}
\begin{enumerate}[label=\arabic*.]
  \item Distinguish a file extension from a file format, with one example.
  \item State the full form of PDF and explain why PDF is called a
    fixed-layout format.
  \item Explain the difference between a text PDF and a scanned PDF, and
    where OCR fits in.
  \item Distinguish encryption from a digital signature in terms of the
    security goal each provides.
  \item What is redaction, and why is drawing a black box over text not a
    safe redaction? What is PDF/A used for?
\end{enumerate}

\subsection*{Solutions}
\begingroup\small
\begin{enumerate}[label=\arabic*.]
  \item The extension is the dot-suffix in the file name, such as
    \texttt{.docx}; it is only a label. The format is the actual structure of
    the data inside the file, such as the DOCX word-processing format. The
    same extension can be renamed without changing the format.
  \item PDF stands for Portable Document Format. It is called a fixed-layout
    format because every item is placed at a fixed position on the page, and
    the fonts and images are carried in the file, so the document looks the
    same on every device and printer.
  \item A text PDF contains real, machine-readable text, so it can be
    selected, copied, and searched. A scanned PDF stores each page as an
    image, so its text cannot be selected or searched. OCR (Optical Character
    Recognition) converts the image of the text into machine-readable text,
    after which the scanned PDF becomes searchable.
  \item Encryption provides confidentiality: it scrambles the file so only
    someone with the correct key or password can read it. A digital signature
    provides authenticity and integrity: it proves who signed the file and
    detects any change made after signing. It does not hide the contents.
  \item Redaction is the permanent removal of sensitive content before a file
    is shared. Drawing a black box over text is not safe because the hidden
    text may still be selectable and copyable underneath. PDF/A is a
    standardised PDF version for long-term archiving that embeds fonts and
    colour information so the document renders the same in the future.
\end{enumerate}
\endgroup
\end{document}
